\documentclass[a4paper,11pt]{article}

\usepackage[utf8]{inputenc}
\usepackage[T1]{fontenc}
\usepackage{amsmath,amssymb}
\usepackage{geometry}
\geometry{margin=2.2cm}

\usepackage{xcolor}
\usepackage{tcolorbox}
\tcbuselibrary{skins,breakable}

\usepackage{booktabs}
\usepackage{longtable}
\usepackage{array}
\usepackage{tikz}
\usepackage{eso-pic}

\usepackage[
    colorlinks,
    citecolor=blue,
    urlcolor=blue,
    linkcolor=blue
]{hyperref}


% ============================================================
% PAGE FRAME
% ============================================================

\AddToShipoutPictureBG{%
  \begin{tikzpicture}[remember picture,overlay]
    \draw[
        black!55,
        line width=0.7pt,
        rounded corners=2mm
    ]
      ([xshift=0.75cm,yshift=0.75cm]current page.south west)
      rectangle
      ([xshift=-0.75cm,yshift=-0.75cm]current page.north east);
  \end{tikzpicture}%
}


% ============================================================
% REVIEWER / EDITOR COMMENT BOX
% ============================================================

\newtcolorbox{reviewerscomment}{
    colback=blue!5,
    colframe=blue!30!yellow,
    arc=3mm,
    left=5mm,
    right=5mm,
    top=3mm,
    bottom=3mm,
    fonttitle=\bfseries,
    title=Reviewer Comment,
    breakable
}


% ============================================================
% RESPONSE BOX
% ============================================================

\newtcolorbox{myresponse}{
    colback=green!5,
    colframe=green!70!black,
    arc=3mm,
    left=5mm,
    right=5mm,
    top=3mm,
    bottom=3mm,
    fonttitle=\bfseries,
    title=Response,
    breakable
}


% ============================================================
% COMMANDS
% ============================================================

% Comment number is supplied manually.
% This prevents global numbering across editors/reviewers.
\newcommand{\reviewercomment}[2]{%
    \begin{reviewerscomment}
        \textbf{Comment #1:} #2
    \end{reviewerscomment}
}

\newcommand{\response}[1]{%
    \begin{myresponse}
        #1
    \end{myresponse}
    \vspace{5mm}
}


% Unnumbered section that still appears in the Contents
\newcommand{\responsesection}[1]{%
    \section*{#1}
    \addcontentsline{toc}{section}{#1}
}


% ============================================================
% DOCUMENT
% ============================================================

\begin{document}

\title{
    Point-by-Point Response to Reviewers\\
    \large Protein-Conditional Normalizing Flows on Manifolds for Backbone Torsion Angle Modeling\\
    \normalsize Machine Learning with Applications
}

\author{Nooshin Akbari Sharak et al.}

\date{\today}

\maketitle

\tableofcontents


% ============================================================
% GENERAL RESPONSE
% ============================================================

\responsesection{General Response}

We sincerely thank the Editor-in-Chief, Associate Editor, and Reviewers for their
careful evaluation of our manuscript and for the constructive comments that
helped us improve both its scientific positioning and empirical evaluation. In
the revised manuscript, we have clarified the scope and novelty of the
protein-conditional formulation, explicitly distinguished our contribution from
the underlying manifold-flow constructions of Rezende et al. (2020), and made
the limitation of identity-based protein embeddings for unseen proteins clear
in the Abstract, Introduction, Method, and Conclusion. We have also strengthened
the empirical evaluation by adding a protein-conditional Riemannian
flow-matching baseline, paired protein-level statistical comparisons, and
real-data density-based clustering analyses on the SCOP datasets. In addition,
we clarified the role of the synthetic experiments, the numerical use of the
grid in the Hellinger-distance calculation, and the reproducibility materials.
All substantive revisions are indicated in colored text in the revised
manuscript. Below, we reproduce the editorial and reviewer comments and provide
our point-by-point responses together with the corresponding locations in the
revised manuscript.


% ============================================================
% SECTION 1-a
% ============================================================

\responsesection{Section 1-a: Complete Comments from the Editor-in-Chief}


\reviewercomment{1}{
I am writing to inform you of the editorial decision regarding your paper
submission MLWA-D-26-00991. to MLWA: Major Revision. If you wish to
submit a revised version of your manuscript, please read this entire editorial
decision letter carefully and address all requirements thoroughly to avoid delays
in the review process.

Your revised submission must include a ``Detailed Response to Reviewers''
document containing the six required sections listed below. Each section must be
clearly labeled and presented in the exact order specified. In addition, please
upload the complete revision package---including the revised manuscript and all
supporting files---through the EM system.

Failure to include any required section or to comply fully with these instructions
may result in rejection of the revised submission.

The six required sections must appear in the following order:

\begin{itemize}
    \item Section 1-a: The two complete comments from the Editor-in-Chief
    \item Section 1-b: The authors' point-by-point responses to the Editor-in-Chief
    \item Section 2-a: The two complete comments from the Associate Editor
    \item Section 2-b: The authors' point-by-point responses to the Associate Editor
    \item Section 3-a: The complete comments from the Reviewers
    \item Section 3-b: The authors' point-by-point responses to the Reviewers
    (organized by individual reviewer)
\end{itemize}

As a verification step, please answer the following question in your response
document:

Does your revised submission fully comply with the required structure outlined
above? Yes $\square$\quad No $\square$

Please also note that a Major Revision decision does not constitute a guarantee
of eventual acceptance, nor does it preclude the possibility of rejection in a
subsequent review round. This decision reflects the reviewers' assessment that
the manuscript has potential but requires significant improvement before further
consideration. Each resubmission will be evaluated independently based on its own
merits, and neither the reviewers nor the journal are bound to a favorable outcome
upon revision.

You are not required to incorporate all references suggested by the reviewers.
If you identify any inappropriate or excessive citation requests, please feel
free to contact me directly.
}


\reviewercomment{2}{
Please confirm (Yes $\square$ No $\square$) that your submission complies with
the following five requirements:

\begin{itemize}

    \item \textbf{Author Information:}
    The EM system and the title page embedded in the full manuscript PDF must
    contain each author's full name, complete institutional affiliation
    (including country), and valid institutional email address. The title page
    should not be submitted separately. All information must be accurate and
    consistent.

    \item \textbf{Institutional Emails:}
    Personal email addresses (e.g., Gmail, Yahoo, Hotmail, or Outlook) are not
    acceptable. The corresponding author is responsible for ensuring complete
    and accurate author information.

    \item \textbf{Authorship:}
    Any authorship changes must be fully justified in the point-by-point response
    and supported by written consent from all authors, including those added or
    removed. Unjustified changes may result in rejection.

    \item \textbf{References:}
    The manuscript must follow the MLWA reference style, based on APA referencing
    style as specified on the MLWA homepage.

    \item \textbf{Manuscript Format:}
    The manuscript must comply with the MLWA Guidelines on Manuscript Format,
    including the required single-column layout.

\end{itemize}

To submit you're a complete and compliant revised manuscript, please log in as
an author at \url{https://www.editorialmanager.com/mlwa/}, and navigate to the
``Submissions Needing Revision'' folder.

To submit you're a complete and compliant revised manuscript, please log in as
an author at \url{https://www.editorialmanager.com/mlwa/}, and navigate to the
``Submissions Needing Revision'' folder.
}


% ============================================================
% SECTION 1-b
% ============================================================

\responsesection{Section 1-b: Authors' Point-by-Point Responses to the Editor-in-Chief}


\response{
\textbf{Response to Comment 1.}

We thank the Editor-in-Chief for clearly specifying the requirements for the
revised submission. We have followed the required structure for the
``Detailed Response to Reviewers'' document and organized it in the exact
sequence requested: Section 1-a contains the two complete comments from the
Editor-in-Chief; Section 1-b contains our point-by-point responses to the
Editor-in-Chief; Section 2-a contains the two complete comments from the
Associate Editor; Section 2-b contains our point-by-point responses to the
Associate Editor; Section 3-a contains the complete comments from the Reviewers;
and Section 3-b contains our point-by-point responses organized separately by
reviewer.

We have also revised the manuscript in response to the editorial and reviewer
comments and will submit the complete revision package, including the revised
manuscript and supporting files, through the Editorial Manager system.

\textbf{Verification:}
Does your revised submission fully comply with the required structure outlined
above?
\textbf{Yes $\boxtimes$\quad No $\square$}
}


\response{
\textbf{Response to Comment 2.}
\textbf{Yes $\boxtimes$\quad No $\square$}

We confirm that the revised submission complies with the five requirements
specified by the Editor-in-Chief. The full names, complete institutional
affiliations including country, and valid institutional e-mail addresses of all
authors have been carefully verified for consistency between the manuscript and
the Editorial Manager system. All authors use institutional e-mail addresses.

No changes have been made to the authorship of the manuscript. Specifically,
the number of authors, author order, and institutional affiliations remain
unchanged from the original submission; therefore, no authorship-change
justification or consent documentation is required.

The manuscript is prepared in the required single-column format. The reference
list has been carefully checked and formatted to conform to the MLWA APA-based
reference style, including alphabetical organization and verification of
bibliographic information and DOI information where available.

\textbf{Location:}
Title page and author information; reference list; manuscript-wide formatting.
}


% ============================================================
% SECTION 2-a
% ============================================================

\responsesection{Section 2-a: Complete Comments from the Associate Editor}


\reviewercomment{1}{
Thank you for submitting your manuscript to MLWA. Following a careful review,
the reviewers have identified several substantive issues that must be adequately
addressed before the manuscript can be considered further. Please revise the
manuscript thoroughly and provide a clear, point-by-point response to each
reviewer comment. All revisions should be clearly indicated using colored text
in the manuscript. Kindly ensure that the track changes function is turned off
prior to submission.
}


\reviewercomment{2}{
Please confirm (Yes $\square$ No $\square$) whether Section \#1-a of your
revised submission includes the complete and unaltered text of two comments
from the Editor-in-Chief. Please note that the complete and verbatim inclusion
of the Editor-in-Chief's comments is a mandatory journal requirement. Any
omission or modification of these comments is considered a breach of editorial
ethics and is not acceptable.
}


% ============================================================
% SECTION 2-b
% ============================================================

\responsesection{Section 2-b: Authors' Point-by-Point Responses to the Associate Editor}


\response{
\textbf{Response to Comment 1.}

We thank the Associate Editor for the careful assessment of our manuscript and
for emphasizing the importance of a thorough revision. We have carefully
considered all comments raised by the reviewers and have undertaken a substantial
revision of the manuscript accordingly.

In particular, we have prepared detailed, point-by-point responses to each
reviewer comment in Section 3-b of this document. The corresponding changes in
the revised manuscript are clearly indicated using colored text so that they can
be readily identified by the Editor and Reviewers. The final revised manuscript
will be submitted with the track-changes function turned off, as requested.

The substantive scientific issues raised by the reviewers, including the scope
and novelty of the protein-conditional formulation, its limitation with respect
to unseen proteins, comparison with a modern generative-modeling baseline,
statistical assessment of model performance, real-data clustering validation,
and reproducibility, are addressed individually in the responses to the
reviewers and through corresponding revisions to the manuscript.

\textbf{Location:}
The revisions are indicated in colored text throughout the revised manuscript,
and detailed locations of individual changes are reported under the corresponding
responses in Section 3-b.
}


\response{
\textbf{Response to Comment 2.}
\textbf{Yes $\boxtimes$\quad No $\square$}

We confirm that Section 1-a of this response document includes the complete and
unaltered text of both comments from the Editor-in-Chief. The comments have been
reproduced verbatim from the editorial decision letter, without omission,
summarization, paraphrasing, or substantive modification.

We have retained the complete wording of the Editor-in-Chief's comments,
including the submission instructions and compliance requirements, to ensure
full adherence to the journal's mandatory response-document structure.

\textbf{Location:}
Section 1-a, ``Complete Comments from the Editor-in-Chief,'' Comments 1 and 2.
}

% ============================================================
% SECTION 3-a
% ============================================================

\responsesection{Section 3-a: Complete Comments from the Reviewers}

\subsection*{Reviewer 1}

\reviewercomment{1}{
The core idea regarding conditioning normalizing flows on learned embeddings
is conceptually straightforward. The heavy theoretical lifting (manifold flows
on circles, tori, spheres) is largely borrowed from Rezende et al. [2020].
The protein-conditional extension is useful but incremental.
}

\reviewercomment{2}{
The embedding-based conditioning mechanism (a simple lookup table) is
lightweight to the point of being limiting. The authors explicitly acknowledge
this cannot generalize to unseen proteins---a major practical constraint that
undermines real-world applicability.
}

\reviewercomment{3}{
In order to strengthen the empirical foundations, the authors should:

1. Addition of at least one modern baseline (e.g., flow matching or
score-based manifold model).
}

\reviewercomment{4}{
In order to strengthen the empirical foundations, the authors should:

2. Explicit discussion of the zero-shot limitation in the
abstract/conclusion, not just Section 2.5.
}

\reviewercomment{5}{
In order to strengthen the empirical foundations, the authors should:

3. Real-data clustering validation on SCOP datasets (currently only synthetic).
}

\reviewercomment{6}{
Considering the scale between 1,591 proteins and $\sim$732K residues
(SCOP Easy) is modest by contemporary standards. AlphaFold2-era structural
biology operates at hundreds of millions of structures.

This work is unlikely to shift practice in protein structure modeling. It
serves more as a principled statistical foundation than a tool practitioners
will immediately adopt.
}

\subsection*{Reviewer 2}

\reviewercomment{1}{
The manuscript presents a technically sound extension of neural circular
spline flows to protein-conditional settings. Still, its novelty is severely
limited because it directly adopts the core manifold flow architecture from
Rezende et al. (2020) with only minor modifications (learned protein
embeddings).
}

\reviewercomment{2}{
The GitHub repository link is non-functional (CFFF), raising reproducibility
concerns.
}

\reviewercomment{3}{
The conditioning mechanism via protein embeddings fails to generalise to
unseen proteins, contradicting the claimed promise of "protein-specific"
modelling.
}

\reviewercomment{4}{
Additionally, the synthetic experiments use overly simplistic distributions
(wrapped normal mixtures, conditional von Mises) that poorly represent the
complex multimodal nature of real protein torsion angles.
}

\reviewercomment{5}{
At the same time, the real-data improvements (NLL differences of 0.05-0.1)
are modest and not statistically validated.
}

\reviewercomment{6}{
The clustering application, while interesting, uses density estimates on
fixed grids rather than the exact likelihood formulation the paper claims
as its advantage.
}

% ============================================================
% SECTION 3-b
% ============================================================

\responsesection{Section 3-b: Authors' Point-by-Point Responses to the Reviewers}

% ============================================================
% REVIEWER 1
% ============================================================

\subsection*{Reviewer 1}
\addcontentsline{toc}{subsection}{Reviewer 1}


% ------------------------------------------------------------
% Reviewer 1 - Comment 1
% ------------------------------------------------------------

\reviewercomment{1}{
The core idea regarding conditioning normalizing flows on learned embeddings
is conceptually straightforward. The heavy theoretical lifting (manifold flows
on circles, tori, spheres) is largely borrowed from Rezende et al. [2020].
The protein-conditional extension is useful but incremental.
}

\response{
We thank the Reviewer for this assessment and agree that the underlying
manifold normalizing-flow machinery is not a novel contribution of the present
work. Specifically, the circular spline transformations on $\mathbb{S}^{1}$,
their extension to product manifolds such as $\mathbb{T}^{D}$, and the
sphere--cylinder construction on $\mathbb{S}^{D}$ build on the manifold-flow
framework of Rezende et al.~(2020). We have revised the manuscript to make this
boundary of novelty explicit and to avoid implying that the underlying
geometric construction itself is introduced here.

The methodological contribution of the present study is instead the
\emph{protein-conditional} formulation built on this framework, in which
learnable protein-specific embeddings are jointly optimized with the flow
parameters. This allows a single shared manifold-aware model to represent
distinct protein-specific torsion-angle distributions while retaining
tractable likelihood evaluation.

We have also strengthened the empirical positioning of this contribution.
In response to the Reviewer's recommendation, we added a direct comparison
with a modern protein-conditional Riemannian flow-matching baseline under a
matched experimental protocol. This additional experiment allows the revised
manuscript to position PC-NCSF empirically relative to a more recent
continuous-time manifold generative-modeling paradigm rather than relying
solely on comparisons with the unconditional NCSF baseline.

\textbf{Location:} Introduction, immediately following the list of main
contributions; Section~3.1.3, ``Comparison with a Modern Flow-Matching
Baseline,'' and Table~3.
}


% ------------------------------------------------------------
% Reviewer 1 - Comment 2
% ------------------------------------------------------------

\reviewercomment{2}{
The embedding-based conditioning mechanism (a simple lookup table) is
lightweight to the point of being limiting. The authors explicitly acknowledge
this cannot generalize to unseen proteins---a major practical constraint that
undermines real-world applicability.
}

\response{
We agree with the Reviewer that the current protein-identity embedding imposes
an important zero-shot limitation. Because the conditioning representation is
a learned lookup embedding associated with proteins observed during training,
the present model cannot directly generate a protein-specific density for a
previously unseen protein without additional fitting. We therefore agree that
this limitation should not be confined to the methodological description.

In the revised manuscript, we have made this limitation explicit in both the
Abstract and the Conclusion and have moderated the discussion of practical
generalizability accordingly. We now distinguish more clearly between
protein-specific density estimation for proteins represented during model
training and zero-shot prediction for unseen proteins, which is not supported
by the current architecture.

We also identify replacing the identity-based lookup embedding with a
conditioning representation derived from transferable protein information,
such as sequence- or structure-derived features, as an important direction for
future work.

\textbf{Location:} Abstract; Section~2.5, ``Conditional Modeling via Learned
Protein Embeddings''; and Section~4, ``Conclusion'' (Limitations).
}


% ------------------------------------------------------------
% Reviewer 1 - Comment 3
% ------------------------------------------------------------

\reviewercomment{3}{
In order to strengthen the empirical foundations, the authors should:

1. Addition of at least one modern baseline (e.g., flow matching or
score-based manifold model).
}

\response{
We agree and have added a modern flow-matching baseline to the revised
manuscript. Specifically, we implemented a Protein-Conditional Riemannian Flow
Matching (PC-RFM) model and evaluated it under the same protein-conditioning
setting and the same train/validation split used for the corresponding
PC-NCSF experiment.

On the paired protein-level comparison, 493 proteins were represented in the
validation outputs of both methods. PC-NCSF achieved a lower mean validation
NLL than PC-RFM for 443 of these 493 proteins (89.9\%). The paired difference
was statistically significant using both the Wilcoxon signed-rank test
($p=3.8\times10^{-69}$) and the paired $t$-test
($p=1.1\times10^{-85}$), with a large standardized paired effect
(Cohen's $d=1.09$).

Importantly, the comparison is performed at the \emph{protein level} rather
than treating individual residues as independent observations. This provides
a more appropriate assessment of the consistency of the performance
difference across proteins and directly addresses the need for statistical
validation of the model comparison.

These results have been added as a new experiment and are reported in the
revised manuscript.

\textbf{Location:} Section~3.1.3, ``Comparison with a Modern Flow-Matching
Baseline,'' and Table~3.
}


% ------------------------------------------------------------
% Reviewer 1 - Comment 4
% ------------------------------------------------------------

\reviewercomment{4}{
In order to strengthen the empirical foundations, the authors should:

2. Explicit discussion of the zero-shot limitation in the
abstract/conclusion, not just Section 2.5.
}

\response{
We agree. As also noted in our response to Comment~2, the inability of the
current lookup-based conditioning mechanism to generalize directly to unseen
proteins is an important limitation of the proposed formulation.

We have therefore added an explicit statement of this limitation to both the
Abstract and the Conclusion. These revisions make clear that the current
protein-conditional formulation provides protein-specific densities for
proteins represented during training but does not constitute a zero-shot model
for previously unseen proteins.

\textbf{Location:} Abstract; Section~4, ``Conclusion'' (Limitations).
}


% ------------------------------------------------------------
% Reviewer 1 - Comment 5
% ------------------------------------------------------------

\reviewercomment{5}{
In order to strengthen the empirical foundations, the authors should:

3. Real-data clustering validation on SCOP datasets (currently only synthetic).
}

\response{
We agree that evaluating the density-based clustering procedure on real
structural data provides a substantially more informative assessment of its
downstream utility. We have therefore extended the Hellinger-distance and
Ward-linkage clustering pipeline, previously demonstrated only in the
controlled synthetic experiments, to all four real SCOP benchmark tiers.

Unlike the near-perfect separation observed in some controlled synthetic
settings, clustering performance on the real SCOP datasets was weak, with
Adjusted Rand Index (ARI) values ranging from 0.001 to 0.097 across the four
tiers. We additionally examined alternative linkage choices and found that
the overall conclusion was unchanged.

We have deliberately retained and reported these results despite their weak
performance because they provide an important qualification of the original
downstream-utility claim. The revised manuscript now makes clear that improved
protein-specific density estimation does not necessarily imply that the
resulting density functions form strongly cluster-separable representations of
complex real protein structural classes.

Accordingly, we have moderated the interpretation of the clustering
experiments. Accurate manifold-aware density estimation is now presented as
the primary contribution of the work, whereas density-based clustering is
treated as a secondary exploratory application. The strong synthetic
clustering results remain useful as a controlled demonstration that the
estimated densities can recover deliberately separated distributional
structure, but they are no longer presented as evidence of comparable
classification performance on real SCOP data.

\textbf{Location:} Section~3.5, ``Real-Data Clustering Validation on SCOP,''
Table~8; and Section~4, ``Conclusion'' (Limitations).
}


% ------------------------------------------------------------
% Reviewer 1 - Comment 6
% ------------------------------------------------------------

\reviewercomment{6}{
Considering the scale between 1,591 proteins and $\sim$732K residues
(SCOP Easy) is modest by contemporary standards. AlphaFold2-era structural
biology operates at hundreds of millions of structures.

This work is unlikely to shift practice in protein structure modeling. It
serves more as a principled statistical foundation than a tool practitioners
will immediately adopt.
}

\response{
We appreciate this perspective and agree that the scale of the datasets
considered in this study is modest relative to contemporary large-scale
protein-structure resources. The objective of the present work, however, is
not large-scale protein structure prediction or foundation-model training.
Rather, it is to investigate accurate likelihood-based density estimation for
manifold-valued protein torsion-angle distributions and to determine whether
protein conditioning improves such density estimates within a shared
geometrically appropriate model.

We have revised the manuscript to make this intended scope more explicit and
to avoid overstating the immediate practical impact of the method. In
particular, the Conclusion now characterizes the proposed framework as a
statistical modeling approach for protein-specific angular distributions,
while recognizing that evaluation on substantially larger structural
collections and the development of transferable conditioning mechanisms will
be necessary before broader practical applicability can be established.

\textbf{Location:} Section~4, ``Conclusion'' (Limitations).
}

% ============================================================
% Reviewer 2
% ============================================================

\subsection*{Reviewer 2}
\addcontentsline{toc}{subsection}{Reviewer 2}


% ------------------------------------------------------------
% Comment 1
% ------------------------------------------------------------

\reviewercomment{1}{
The manuscript presents a technically sound extension of neural circular
spline flows to protein-conditional settings. Still, its novelty is severely
limited because it directly adopts the core manifold flow architecture from
Rezende et al. (2020) with only minor modifications (learned protein
embeddings).
}

\response{
We thank the Reviewer for this important comment. We agree that the underlying
manifold normalizing-flow constructions are not a novel contribution of the
present work. In particular, the circular, toroidal, and spherical geometric
constructions used by PC-NCSF build on the framework of Rezende et al.
(2020). We have therefore revised the Introduction to make this boundary of
novelty explicit and to avoid attributing the underlying manifold-flow
architecture itself to the present study.

Our contribution is more specifically the protein-conditional formulation:
learnable protein-identity embeddings are jointly optimized with a shared
manifold flow, allowing a single model to yield distinct, individually
queryable protein-specific angular density estimates while sharing statistical
information across proteins. We have revised the manuscript to state this
contribution more precisely and have moderated the corresponding novelty
claims.

To strengthen the empirical assessment of whether this conditional
formulation provides value beyond the underlying manifold-flow construction,
we have also added a new comparison with Protein-Conditional Riemannian Flow
Matching (PC-RFM), a more recent continuous-time manifold generative-modeling
approach. Under the matched experimental setting, PC-NCSF achieved lower
mean validation NLL for 443 of the 493 proteins (89.9\%), with the paired
difference supported by both the Wilcoxon signed-rank test
($p = 3.8\times10^{-69}$) and the paired $t$-test
($p = 1.1\times10^{-85}$), and a large paired effect size
(Cohen's $d = 1.09$). These results are now reported in the new
Section~3.1.3.

Accordingly, we do not position PC-NCSF as a new manifold-flow geometry, but
rather as a protein-conditional adaptation and empirical evaluation of
established manifold-flow machinery for protein-specific angular density
estimation.

\textbf{Location:} Introduction, immediately following the list of main
contributions; Section~3.1.3.
}

% ------------------------------------------------------------
% Comment 2
% ------------------------------------------------------------

\reviewercomment{2}{
The GitHub repository link is non-functional (CFFF), raising reproducibility
concerns.
}

\response{
We thank the Reviewer for identifying this issue. We have corrected the
repository accessibility problem and verified that the repository is now
publicly accessible at
\url{https://github.com/mamintoosi-papers-codes/CFFF}.

We have also reviewed the repository for consistency with the revised
manuscript. In particular, an internal naming inconsistency in which some
training scripts referred to the conditional model as ``PC-FFF'' has been
corrected to ``PC-NCSF'', and the repository documentation and citation
metadata have been updated accordingly. These revisions are intended to make
the computational materials easier to identify and use alongside the
manuscript and to improve the reproducibility of the reported experiments.

\textbf{Location:} Abstract (implementation and code availability statement);
GitHub repository and repository documentation.
}

% ------------------------------------------------------------
% Comment 3
% ------------------------------------------------------------

\reviewercomment{3}{
The conditioning mechanism via protein embeddings fails to generalise to
unseen proteins, contradicting the claimed promise of "protein-specific"
modelling.
}

\response{
We thank the Reviewer for highlighting this important distinction. We agree
that the current protein-identity embedding is a learned lookup representation
and therefore does not support zero-shot generalization to proteins that were
absent during model training. We have revised the manuscript to state this
limitation explicitly and to clarify the intended scope of the term
``protein-specific.''

In the present formulation, ``protein-specific'' refers to the model's ability
to learn and provide individually queryable conditional density estimates for
each protein represented during training, while sharing the parameters of the
underlying manifold flow across proteins. It is not intended to imply that
PC-NCSF can infer the conformational density of a previously unseen protein
from its sequence or structure.

To prevent this distinction from being overstated, we now explicitly identify
the lack of zero-shot generalization as a limitation in the Abstract,
Introduction, and Conclusion. We also identify replacing the identity-based
lookup embedding with sequence- or structure-derived conditioning
representations as an important direction for future work.

\textbf{Location:} Abstract; Introduction, immediately following the
clarification of the contribution and novelty; Section~2.5; and Conclusion
(Limitations).
}


% ------------------------------------------------------------
% Comment 4
% ------------------------------------------------------------

\reviewercomment{4}{
Additionally, the synthetic experiments use overly simplistic distributions
(wrapped normal mixtures, conditional von Mises) that poorly represent the
complex multimodal nature of real protein torsion angles.
}

\response{
We thank the Reviewer for this important comment. We agree that controlled
synthetic distributions cannot reproduce every aspect of the heterogeneity
and multimodal complexity observed in empirical protein torsion-angle data.
However, we would like to clarify that the synthetic settings were not chosen
as arbitrary simplified distributions. In particular, the first simulation
setting is motivated by the characteristic bimodal angular structure of
consecutive $\mathrm{C}_\alpha$ atoms reported by Hamelryck et al. (2006),
as already described in Section~3.2.1.

The purpose of the simulation study is to provide protein-relevant angular
settings in which the generating densities and group structure are known.
This makes it possible to directly evaluate recovery of the protein-specific
densities and to determine whether density-based clustering recovers the
known ground-truth groups---assessments that cannot be performed in the same
way on empirical protein data because the underlying densities are unknown.

To make this rationale clearer, we have revised the opening paragraph of
Section~3.2 to distinguish the role of the synthetic experiments from that
of the real-data analyses. The synthetic experiments are intended as
controlled methodological evaluations that complement, rather than replace,
the experiments on real protein data in Section~3.3 and the real-data
density-based clustering analysis in Section~3.5.

\textbf{Location:} Section~3.2 (opening paragraph).
}
% ------------------------------------------------------------
% Comment 5
% ------------------------------------------------------------

\reviewercomment{5}{
At the same time, the real-data improvements (NLL differences of 0.05-0.1)
are modest and not statistically validated.
}

\response{
We thank the Reviewer for this important observation. We agree that the
magnitude of the NLL improvement should be distinguished from its statistical
support. To address this concern, we performed a paired protein-level analysis
comparing the unconditional model with the conditional model using
$d_{\mathrm{emb}}=16$.

Specifically, we computed validation NLL separately for each protein and
performed paired comparisons across the 493 proteins represented in both
validation outputs. The conditional model achieved lower NLL for 283 of
493 proteins (57.4\%), with a mean paired improvement (unconditional minus
conditional NLL) of 0.060 nats. The paired difference was statistically
significant under both the Wilcoxon signed-rank test
($p = 5.3\times10^{-9}$) and the paired $t$-test
($p = 2.3\times10^{-12}$), with a small standardized paired effect
(Cohen's $d = 0.32$).

We have added these results to Section~3.1.2. We have also revised the
interpretation to distinguish statistical evidence for a systematic
improvement from its modest effect magnitude, rather than relying solely on
the aggregate validation NLL.

\textbf{Location:} Section~3.1.2, immediately following the comparison of
embedding dimensions with the unconditional baseline.
}


% ------------------------------------------------------------
% Comment 6
% ------------------------------------------------------------

\reviewercomment{6}{
The clustering application, while interesting, uses density estimates on
fixed grids rather than the exact likelihood formulation the paper claims
as its advantage.
}

\response{
We thank the Reviewer for raising this point. We agree that the original
description could give the impression that the learned densities themselves
were approximated on a fixed grid. This is not the case.

The fixed grid is used only to numerically approximate the Hellinger-distance
integral required for the downstream clustering analysis. At each grid point,
the protein-specific density is evaluated directly from the trained PC-NCSF
model using its invertible change-of-variables formulation; the grid is not
used to estimate or approximate the underlying density function. Thus, the
numerical approximation concerns the integration required to obtain pairwise
Hellinger distances, rather than the pointwise density or likelihood
evaluation provided by PC-NCSF.

We have revised Section~3.2.3 to make this distinction explicit and to remove
the potentially ambiguous wording that the estimated densities are
``evaluated on a uniform grid.''

\textbf{Location:} Section~3.2.3, Density-Based Clustering
(Clustering procedure).
}

\end{document}